% Осень26!A17; нумерация внутри столбца без пустых строк.
% Источник: лекция 6, §6 (строка 347)
\ticket{15}{Следствие о сохранении области для отображений в пространство меньшей размерности}
\label{ticket:15}

\begin{theorem}
Пусть $O\subset\mathbb R^m$ открыто, $F:O\to\mathbb R^\ell$ принадлежит классу $C^1$, $\ell\le m$ и $\operatorname{rank}F'(x)=\ell$ для всех $x\in O$.
Тогда $F(O)$ открыто в $\mathbb R^\ell$.
\end{theorem}

\begin{proof}
Для $a\in O$ выберем ненулевой минор порядка $\ell$ в $F'(a)$.
Перенумеровав координаты, считаем, что он расположен в первых $\ell$ столбцах.
По непрерывности производных этот минор не равен нулю в некоторой окрестности $U$ точки $a$.
Определим
\[
  \widetilde F(x)=(F_1(x),\ldots,F_\ell(x),x_{\ell+1},\ldots,x_m).
\]
Матрица $\widetilde F'$ блочно-треугольная; её определитель равен выбранному минору и отличен от нуля на $U$.
По \hyperref[ticket:14]{теореме о сохранении области} множество $\widetilde F(U)$ открыто в $\mathbb R^m$.
Координатная проекция $\pi:\mathbb R^m\to\mathbb R^\ell$ переводит открытые множества в открытые, поскольку $\pi(B(x,r))=B(\pi x,r)$.
Следовательно, $F(U)=\pi(\widetilde F(U))$ --- открытая окрестность $F(a)$, содержащаяся в $F(O)$.
\end{proof}
